\chapter{Descripción del componente desarrollado}

\section{Marco teórico}

\subsection{Definición de un plasma}

\noindent Un plasma consiste en una colección de partículas cargadas (iones y electrones), además de átomos y moléculas neutras en ciertos casos \cite{Umran, Chen, Parks}. La formación de un plasma se produce cuando aumenta la temperatura de un gas neutro, o cuando se somete a campos eléctricos muy grandes, hasta que sus átomos y moléculas se ionizan \cite{Umran, Chen}. Por lo anterior, un plasma también puede considerarse como un gas ionizado. Por otro lado, la naturaleza cargada de las partículas de un plasma hace que sus propiedades sean diferentes a un gas normal, un líquido o un sólido, razón por la que se conoce como el cuarto estado de la materia \cite{Umran}. Es más, la mayoría de la materia ordinaria del universo se encuentra en estado de plasma \cite{Chen}.

% Un plasma tiene ciertas características fundamentales que lo diferencian de un gas ionizado. Las partículas cargadas de un plasma tienden a apantallar los campos eléctricos creados por cualquier perturbación eléctrica dentro de una distancia mucho menor que la dimensión característica del plasma \cite{Umran, Chen, Parks}. La distancia a la que sucede este apantallamiento se conoce como longitud de Debye. A distancias mayores que la longitud de Debye las interacciones individuales entre las partículas cargadas dejan de ser dominantes y el plasma tiene un comportamiento “colectivo” \cite{Kamide}. Esto quiere decir que el movimiento de las partículas no depende de las condiciones locales sino del plasma en su totalidad. Además, debido al apantallamiento las cargas se reorganizan de tal manera que en un plasma la densidad de iones positivos es aproximadamente igual a la densidad de electrones, y se dice que un plasma es un gas cuasi neutral \cite{Parks, Umran}. Esta tendencia a mantener la neutralidad provoca que los electrones de un plasma oscilen a una frecuencia característica en la que ocurren las perturbaciones electrostáticas en los plasmas.  Esta frecuencia de oscilación de electrones es mayor que la frecuencia a la que colisionan las partículas en el plasma, pues de lo contrario se pierde el comportamiento colectivo \cite{Chen}.


\subsection{Descripciones de un plasma}

\noindent En física de plasmas existen varios modelos para describir a un plasma. Dependiendo de las condiciones supuestas y de las diferentes aproximaciones que se realicen, un plasma se puede estudiar a través del movimiento de partículas individuales sujetas a la fuerza de Lorentz, mediante funciones de distribución a partir de la ecuación de Boltzmann debido al gran número de partículas presentes en el plasma, o por medio de las ecuaciones de la mecánica de fluidos, en conjunto con las ecuaciones del electromagnetismo, enfocándose en las propiedades colectivas del plasma \cite{Umran, Chen, Krall}.   

El enfoque fluido es muy importante al proporcionar una descripción \textit{macroscópica} de un plasma. Esta descripción es útil por las propiedades comunes entre un plasma y un fluido conductor, mostrando movimientos colectivos como, por ejemplo, ondas \cite{Krall}. En este modelo se considera a todo el plasma como un fluido de un solo tipo, enfocándose en las propiedades colectivas y despreciando la diferencia entre las partículas cargadas \cite{Umran, Chen}. Este enfoque se conoce como magnetohidrodinámica y es muy útil para describir un amplio rango de comportamientos de un plasma, abarcando la física espacial y de fusión \cite{Umran}. 

\subsection{El Sol y el Viento Solar}

\noindent El Sol es la principal fuente de energía que mantiene la vida en el planeta Tierra.  Es un plasma conformado con alrededor de $90\%$ de hidrógeno, $10\%$ de helio y una pequeña proporción de otros elementos tales como carbono, nitrógeno, y oxígeno \cite{Parks}. La temperatura en su interior ronda los $15 \times 10^{6} K$, permitiéndole producir energía a través de la fusión nuclear de hidrógeno en helio \cite{Cravens}. Una pequeña proporción de la energía generada viaja alrededor de $1.496 \times 10^{8} km$ para llegar desde el Sol hasta la Tierra \cite{Kamide}. Esta distancia define la unidad astronómica (UA).

La atmósfera del Sol se encuentra dividida en la fotosfera, la cromosfera, y la corona \cite{Parks, Cravens}. Esta última capa se sitúa sobre la otras dos y se extiende hacia el espacio interplanetario \cite{Kamide}. La expansión de la corona produce un flujo de partículas cargadas que son expulsadas por el Sol, comúnmente conocido como viento solar, y se descubrió en los comienzos de la era espacial \cite{Parks}. El viento solar es una consecuencia natural de estrellas calientes como el Sol, pues cuando la temperatura de la corona es de alrededor de $10^6 K$, casi el $50\%$ de los electrones tienen velocidades térmicas que exceden la velocidad de escape gravitacional del Sol, aunque en el caso de los protones es de menos del $1\%$ \cite{Parks, encrenaz_solar_2004}, debido a su mayor masa. Sin embargo, el exceso de carga positiva se contrarresta por la generación de un enorme campo eléctrico que provoca la aceleración de protones, expulsándolos del Sol \cite{encrenaz_solar_2004}. Por tanto, la expansión de la corona es una manera en la que el Sol mantiene su neutralidad de carga \cite{Parks}. 

El Sol tiene un campo magnético que a grandes escalas se aproxima a un dipolo magnético \cite{Cravens}, cuya intensidad superficial promedio es de alrededor de $10^{-4} T$ \cite{Parks, Kamide, Cravens}. En su expansión a través del espacio interplanetario, el viento solar lleva consigo el campo magnético del Sol que se extiende radialmente hacia afuera \cite{Kamide} y cubre todo el sistema solar, protegiendo a los planetas de las partículas de alta energía del espacio interestelar. El límite en donde se encuentra el viento solar con el viento interestelar se llama heliopausa, mientras que la región en su interior se denomina heliosfera \cite{Parks, encrenaz_solar_2004}, debido a la mayor influencia del Sol. Mediciones realizadas por magnetómetros en sondas espaciales muestran que la intensidad de campo magnético interplanetario alrededor de la Tierra es de unos pocos $10^{-9} T$ \cite{Parks, Kamide}. 

\subsection{Misiones de Exploración Espacial y Series Temporales}

\noindent Las sondas espaciales permiten mejorar la comprensión del espacio a través de la recopilación de información de gran interés científico. Los instrumentos a bordo son capaces de realizar mediciones y observaciones que en la Tierra no serían posibles. En particular, las misiones Parker Solar Probe (PSP), a cargo de la NASA y lanzada en 2018 \cite{fox_solar_2016}, y Solar Orbiter (SolO), dirigida por la ESA y lanzada en 2020 \cite{muller_solar_2020}, están encargadas de estudiar la heliosfera interior, enfocándose en el Sol. Estas sondas espaciales llevan instrumentos de detección a distancia (por ejemplo, telescopios) e in-situ (por ejemplo, magnetómetros), para realizar observaciones de la corona solar y para medir el viento solar, las partículas energéticas y los campos electromagnéticos \cite{fox_solar_2016, muller_solar_2020}, respectivamente. En conjunto, las sondas PSP y SolO posibilitan profundizar el entendimiento sobre la atmósfera del Sol y su actividad \cite{velli_understanding_2020}. 
    
Los instrumentos de las sondas espaciales registran los valores de diferentes magnitudes físicas a intervalos de tiempo iguales. A este tipo de mediciones se les conoce como series temporales \cite{Priestley, chatfield_analysis_2019}. De las series temporales registradas por las sondas espaciales son de especial interés las mediciones del campo magnético. 

La dinámica de los plasmas espaciales es un proceso intrínsecamente no lineal \cite{Kamide, bruno_turbulence_2016, matthaeus_intermittency_2015} y en general, no estacionario \cite{chatfield_analysis_2019, Priestley, alberti_multifractal_2019}. Esto último significa que las propiedades estadísticas (por ejemplo, la media y la varianza) de las series temporales cambian con el tiempo \cite{chatfield_analysis_2019, Priestley} y, por consiguiente, los datos no son estacionarios. Por otro lado, se dice que un sistema lineal tiene ``datos lineales”.          

\subsection{Turbulencia}

\noindent El movimiento de los fluidos se puede dividir en dos tipos, laminares o turbulentos. El primer caso implica un movimiento regular y suave del fluido, es decir, un movimiento ``predecible'' \cite{Landau, Cengel, davidson_turbulence_2015}. Por otro lado, un flujo turbulento se caracteriza por tener un carácter aleatorio e impredecible \cite{Landau, Cengel, davidson_turbulence_2015}. Por ejemplo, el flujo de humo de un cigarro que se deja consumir es laminar al principio, pero a medida que el humo se desplaza se vuelve turbulento formando remolinos de diversos tamaños que se mueven aleatoriamente \cite{Cengel}. 

La turbulencia es un fenómeno que se puede manifestar en cualquier fluido, pues el movimiento de un fluido es casi siempre inestable \cite{davidson_turbulence_2015}. No existe una definición rigurosa sobre la turbulencia, pero se distinguen características que permiten cierto entendimiento. El carácter impredecible de la turbulencia se debe a la sensibilidad a las condiciones iniciales, lo cual refleja su naturaleza caótica \cite{davidson_turbulence_2015}. De igual manera, grandes fluctuaciones en la velocidad del fluido son propias de la turbulencia \cite{davidson_turbulence_2015, Cengel}. Además, es común la presencia de remolinos que se forman y desplazan de manera aleatoria, como se mencionó en el ejemplo anterior. Finalmente, en un flujo turbulento, como resultado de la difusión del fluido, se produce disipación de energía, desde remolinos grandes hacia remolinos cada vez más pequeños, en la denominada cascada de energía de Richardson \cite{davidson_turbulence_2015, bruno_turbulence_2016}. 

\subsubsection{Turbulencia en el Viento Solar}

\noindent Los plasmas espaciales y los fluidos atmosféricos son sistemas turbulentos en evolución dinámica, cuyas interacciones son no lineales \cite{Kamide}. La naturaleza de las interacciones en estos sistemas es evidente cuando se tiene presente que algunas de las ecuaciones de la mecánica de fluidos son ecuaciones diferenciales no lineales parciales acopladas. En el caso del viento solar, las mediciones realizadas por sondas espaciales han demostrado que la turbulencia es un fenómeno que muestra la presencia de fluctuaciones a pequeña escala (espaciales y temporales) en el campo magnético, el campo de velocidades y la densidad de protones \cite{bruno_turbulence_2016, sorriso-valvo_intermittency_1999, hnat_intermittency_2003}. En el viento solar, la turbulencia se desarrolla fuertemente durante el desplazamiento y expansión a través del espacio interplanetario \cite{coleman_turbulence_1968, bruno_solar_2013, bruno_turbulence_2016}, similar al ejemplo del humo de cigarro.

\section{Motivación}

\noindent El alineamiento de sondas espaciales es un evento poco común del cual solamente se tienen unos pocos casos registrados \cite{damicis_radial_2010, bruno_spectral_2014, bruno_radial_2014, schwartz_radial_1983, zank_theory_2017, telloni_evolution_2021} desde que se empezaron a realizar mediciones in situ de la actividad magnética del Sol hace unas pocas décadas. Estos sucesos son de gran importancia para la investigación de las características de los plasmas espaciales porque permiten estudiar la evolución del plasma a través de la heliosfera.

En particular, el seguimiento de un plasma espacial permite estudiar las propiedades de la turbulencia del viento solar, lo cual es crucial para el entendimiento de los plasmas y en general, de los fluidos. Esto se realiza analizando la influencia de la turbulencia sobre las variables que se registran en series temporales. Recientemente, las sondas Parker Solar Probe y Solar Orbiter, tuvieron su primer alineamiento radial a una distancia de $ \sim 0.1 UA$ y $ \sim 1 UA$ \cite{telloni_evolution_2021}, respectivamente. Las mediciones del campo magnético del viento solar mostraron que el plasma pasa de un estado poco turbulento cerca del Sol hacia un comportamiento con una turbulencia altamente desarrollada e intermitente en las cercanías de la Tierra \cite{telloni_evolution_2021}. La motivación detrás de este trabajo es investigar la evolución de la parte no turbulenta del campo magnético cuando el viento solar se desplaza radialmente desde el Sol hasta la Tierra. Esta parte no turbulenta se puede cuantificar mediante un filtrado de las altas frecuencias de las series temporales asociadas.

El tipo de análisis que se realiza sobre las series temporales es crucial en cuanto a las características que se pretenden estudiar. La correlación es una medida estadística de uso frecuente para analizar series temporales. Esta mide el grado de similitud entre dos series temporales \cite{chatfield_analysis_2019}. Por ejemplo, en el seguimiento de un plasma espacial, las series temporales de cada sonda deben estar altamente correlacionadas para concluir que pertenecen al mismo plasma.

Es común que el análisis de las series temporales de los plasmas espaciales se realice mediante métodos estándar como aquellos basados en la transformada de Fourier o transformadas wavelet \cite{chatfield_analysis_2019}. Sin embargo, los primeros solamente son útiles para datos estacionarios \cite{chatfield_analysis_2019}, mientras que los últimos están restringidos a datos con escalas temporales definidas \cite{huang_empirical_1998, Percival}. Además, algunos de los métodos estándar son estrictamente útiles para datos que provienen de sistemas lineales, o son muy especializados para explorar nuevas características en los datos. Por otro lado, las series temporales de las magnitudes físicas del viento solar son no lineales, no estacionarias y dependen de la escala \cite{de_michelis_multi-scale_2012, carbone_arbitrary-order_2018, alberti_multifractal_2019, carbone_statistical_2020, carbone_statistical_2021}, es decir, de los intervalos espaciales y temporales en donde ocurren ciertos procesos físicos.   

El método Empirical Mode Decomposition (EMD) puede ser útil para el análisis de series temporales magnéticas \cite{de_michelis_multi-scale_2012, carbone_arbitrary-order_2018, alberti_multifractal_2019, carbone_statistical_2020, carbone_statistical_2021}. Este método descompone una serie temporal en una suma de componentes oscilatorias denominadas funciones de modo intrínseco (IMFs), que permiten una buena reconstrucción de la serie. Reconstrucciones parciales (no todas las IMFs) hacen posible filtrar las componentes de alta frecuencia de las series temporales. Debido al carácter empírico del método, las características locales de la escala de tiempo de los datos son consideradas \cite{huang_empirical_1998}, y esto hace posible estudiar un mayor tipo de datos. Por esta razón, en este trabajo se utiliza el EMD para analizar las series temporales del viento solar durante el alineamiento de las sondas espaciales PSP y SolO. 

\section{Objetivos}

\subsection{Objetivo General}

\noindent En este estudio se busca identificar posibles correlaciones relevantes entre las reconstrucciones a baja frecuencia de las series temporales de campo magnético del viento solar, en los alrededores del Sol y la Tierra (medidas por PSP y SolO, respectivamente), obtenidas mediante las IMFs de la descomposición del método Empirical Mode Decomposition. De esta manera se quiere estimar si se puede disminuir la aparente aleatoriedad que aporta la turbulencia sobre las series temporales y conservar resultados significativos.

%Para cumplir el objetivo general de este trabajo se plantean los siguientes objetivos específicos:

\subsection{Objetivos Específicos}

\begin{enumerate}
    \item Implementar y aplicar el método Empirical Mode Decomposition a las series temporales.
    \item Implementar el método Cross-Correlation y aplicarlo entre reconstrucciones de las series temporales realizadas a partir de sus IMFs de menor frecuencia.  
\end{enumerate}

\section{Alcance}

\noindent El presente trabajo pretende ser un estudio exploratorio de la evolución de la parte no turbulenta de las series temporales del campo magnético al filtrar las componentes de alta frecuencia mediante el EMD. Además, se indaga en el número de IMFs necesario para reconstruir las series temporales y obtener resultados significativos en el cálculo de las correlaciones entre las series.

En este trabajo se utilizaron datos que derivan de otro estudio de series temporales de campo magnético, que también corresponden con el alineamiento de las sondas espaciales PSP y SolO \cite{telloni_evolution_2021}. Los datos recibidos fueron preprocesados por el grupo de investigación anterior por lo que se pudo omitir este paso en el análisis de las series temporales. Esto permitió continuar con la implementación y aplicación del método EMD, que en este caso se realizó en Fortran. También se utilizaron las librerías de Matlab y Python para comparar y validar la implementación realizada. 

El método Cross-Correlation es uno de los métodos estándar más utilizados para correlacionar series temporales \cite{chatfield_analysis_2019}. Este método se implementó en Python debido a la eficiencia de tiempo, manejo y almacenamiento de variables, para el caso particular del cálculo de las correlaciones propuestas. 